% Conclusion. Thesis owned. Readable on its own. One short paragraph per
% research question of the introduction, then the proposal, then stop. No new
% material, no closing offer.

A task fixes what a network computes and leaves open how much of the network
does the computing. With the target function, the data, the architecture and
the width all held fixed, the rule that assigns credit during training settles
how many neurons carry a grokked solution, and the settlement is not small:
the same Fourier construction is implemented by about $154$ of $256$ neurons
or by about $254$, depending on the rule.

Which rules move it can be read off the form of their update. A neuron drops
out of use when the credit it receives is small relative to the part of its
own update that does not come from the gradient, and that one condition covers
three mechanisms which had no common account before. Only a normalising factor
shared by exactly one neuron can change that ratio, which is a yes or no about
a rule rather than a ranking among rules, and it places the optimizers in
general use on one side of the line and two rules on the other.

The choice is not free. A computation spread over more neurons tolerates the
deletion of neurons better. A computation held on fewer prunes further and
adapts about twice as fast to a task that needs new combinations of the
features already learned, though not to one that only changes which answer
each feature votes for. And moving from the concentrated circuit to the spread
one costs about two points of test accuracy on modular addition, from $0.998$
to $0.976$.

The learning rule is not the only component that leaves structure behind. With
backpropagation held fixed and the training objective varied, the set of token
pairs the objective admits settles whether a Transformer's attention weights
stay close to symmetric or whether a few columns come to dominate, in models
at the scale the field works at.

A measurement taken of a trained network is only as good as the readout behind
it. The probe used to score latent world models responds largely to the norm
of a predicted latent, a quantity a cosine training loss never fixes, and
rescaling each prediction to a reference norm taken from the encodings removes
most of the error without retraining anything.

That standard applies to this thesis's own measure, and applying it changes
the answer in part. Measures tied to the computation place the untouched
network lower than the structural measure does while moving in the same
direction under every intervention, so the direction of the central result
survives the check and the level reported by the structural measure is the
largest of the three.

The last chapter is a proposal rather than a result. An architecture in which
every active neuron is held in place by several others filters noise that
nothing supports, completes figures that are partly hidden, and recombines
fragments learned separately into patterns it was never trained on, all
without a learning rule being tuned to produce that structure. It does not
scale beyond line drawings, its features are designed rather than learned, and
it defines no dynamics over its representation.

What the thesis establishes is that how far a computation spreads over a
network is a property of how the network was trained, that the property can be
chosen, and that choosing it has a price which can be measured and stated.
